% Pursuing Stacks (cotes 134-2 to 134-8): archive cote 134-6, site page 92 (PDF page 93), typescript page 410.
% typescript. Draft compared: ps5.tex:6183-6237.
% AI-transcribed by Claude Opus 5.5 (claude-opus-5-5, adaptive thinking, effort high), Message Batches API, batch ps3 (msgbatch_01XdFmuVVBAcgc7EFxdtCqsj), 2026-11-12; reviewed against both supplied scan halves; see CHECK_DECISIONS.tsv.
\scanpage{134-6/92}
\mspage{410}

operation is it understood.  I'll have to play around a little more with
it for being really at ease.  It's worth the while,  as \deleted{\xout{(65)}} the \uline{Hom}
and corresponding Hom  operation is the key operation for expressing
cohomology \inserted[typed]{of $A^\wedge$} (with coefficients in K),  whereas the \hand{$\mathbin{\tikz[baseline=-.5ex]{\node[circle,draw,inner sep=.04em] {\ensuremath{\textcolor{black}{*}}};}}$} and corresponding
$*$  operation is the key for expressing homology \inserted[typed]{of $A^\wedge$} (with coefficients in M ,
where K and M are the sheaves occuring in (65) and (67) respectively).

A typical special case is the one when A is the final category,
hence B=A and $A^\wedge_{\hand{\mathcal{M}}}\hand{\simeq \mathcal{M}}\deleted{\illegible{1}}$, in which case (65) and (67) are the two familiar
exterior operations of (Ab) on any additive category stable under the
two types of small limits
\[ (L,\ X)\hand{\mapsto}\deleted{\xout{\illegible{1}_{\underline{Z}}}}\ \underline{\mathrm{Hom}}_{\underline{Z}}(L,X) \quad : \quad (Ab)^{o}\ \mathrm{x}\ \hand{\mathcal{M}}\hand{\to}\ \hand{\mathcal{M}}, \]
and
\[ (X,L)\hand{\mapsto} X\otimes_{\underline{Z}}L \quad : \quad (Ab)\ \mathrm{x}\ \hand{\mathcal{M}}\hand{\to}\hand{\mathcal{M}}, \]
each one of these two operations being deducible from the other by the
usual device of replacing \hand{$\mathcal{M}$} by the dual category \hand{$\mathcal{M}$}$^{o}$ .  When \hand{$\mathcal{M}$}= (Ab) ,
these are just the usual inter\unsure{b|h|n}al  Hom = \uline{Hom}  and tensor product opera-%
tions.  This very particular case shows at once that we should'nt expect
in general the operations  $\underline{\mathrm{Hom}}_{\underline{Z}}(L,-)$ of $A^\wedge_{ab}$  and  $-*L'$  of $B^\wedge_{ab}$  upon
$A^\wedge_{\hand{\mathcal{M}}}$  to commute up to isomorphism  -  we should'nt expect,  for a given
L \inserted[typed]{in $A^\wedge_{ab}$} or a given L' in $B^\wedge_{ab}$ the commutation relation to hold,  except when
this object is ``projective of finite presentation''  i.e. is a direct fac-%
tor of an object of Add(A)  resp. \inserted{of} Add(B) \hand{$\approxeq$} Add(A)$^{o}$ .  Another fact beco-%
ming evident by this particu\unsure{b|l}ar case,  is that whereas it is true that
in the equivalence (64)
\[ B^\wedge_{ab}\ \hand{\approxeq}\ \underline{\mathrm{Hom}}_{\hand{!}}(A^\wedge_{ab},\ (Ab)) \]
\inserted[typed]{a} projective object \inserted[typed]{L'}\xout{s} in $B^\wedge_{ab}$ gives rise to \inserted[typed]{a} functor\xout{s}  $A^\wedge_{ab}\hand{\longrightarrow}$ (Ab) which is
\hand{\uline{\textcolor{black}{exact}}} (besides commuting to \deleted{any} small direct limits)  -  and even to a
functor commuting to small inverse limits if L is ultraprojective i.e.
projective and of finite presentation -  the converse to this (as con-%
templated on page 381) does \hand{\uline{\textcolor{black}{not}}} hold true.  Indeed,  in the particular
case A =\hand{$\Delta$}$_{o}$ ,  when L' is just an \xout{abelian group} object in (Ab),  the
exactness property e\over{n}{n}visioned i.e. exactness of the functor  $M\hand{\mapsto}\over{M}{M}\otimes_{\underline{Z}}L'$
from (Ab) to itself,  only means that L' is a \hand{\uline{\textcolor{black}{flat}}}  $\underline{Z}$-module (i.e.
torsion-free),  which does \hand{\uline{\textcolor{black}{not}}} imply that it is projective (i.e. free).

\margin{92}
